\documentclass[11pt, a4paper]{article}
\usepackage[utf8]{inputenc}
\usepackage[T1]{fontenc}
\usepackage[spanish]{babel}
\usepackage{amsmath, amssymb}
\usepackage{geometry}
\usepackage{enumitem}
\usepackage{titlesec}
\usepackage{parskip}

\geometry{margin=2.5cm}

\titleformat{\section}{\large\bfseries}{}{0em}{}
\titleformat{\subsection}{\normalsize\bfseries}{}{0em}{}
\titleformat{\subsubsection}{\normalsize\itshape}{}{0em}{}

\newcommand{\argmax}{\operatorname*{arg\,m\acute{a}x}}

\title{\textbf{Resumen -- Aprendizaje por Refuerzo} \\
\large Inteligencia Artificial y Neurociencias \\ Universidad Torcuato Di Tella}
\author{}
\date{}

\begin{document}
\maketitle

%% ============================================================
\section{Clase 0 -- Introducci\'on al Aprendizaje por Refuerzo}
%% ============================================================

El \textbf{Aprendizaje por Refuerzo} (RL) es un paradigma en el que un \textbf{agente} aprende a comportarse interactuando con un \textbf{ambiente}: ejecuta \textbf{acciones}, observa el nuevo estado del ambiente y recibe \textbf{recompensas}. El objetivo es aprender a actuar de modo de maximizar la recompensa acumulada.

\subsection{Tres caracter\'isticas distintivas}
\begin{itemize}[nosep]
    \item \textbf{Prueba y error:} nadie le dice al agente cu\'al es la acci\'on correcta; la descubre probando y midiendo consecuencias.
    \item \textbf{Recompensas demoradas:} una acci\'on puede tener efecto sobre recompensas futuras, no solo la inmediata (el problema de \emph{asignaci\'on de cr\'edito}).
    \item \textbf{Balance exploraci\'on/explotaci\'on:} el agente debe \emph{explotar} lo que ya sabe que da buenas recompensas, pero tambi\'en \emph{explorar} para descubrir mejores opciones.
\end{itemize}

\subsection{RL vs.\ otros paradigmas}
\begin{itemize}[nosep]
    \item \textbf{Aprendizaje supervisado:} aprende de ejemplos etiquetados (entrada $\to$ salida correcta). En RL no hay etiquetas, solo se\~nales de recompensa.
    \item \textbf{Aprendizaje no supervisado:} busca estructura oculta en datos sin etiquetar. RL no busca estructura, busca \emph{maximizar recompensa}.
\end{itemize}
Bibliograf\'ia de referencia: Sutton \& Barto, \textit{Reinforcement Learning: An Introduction} (2018).

%% ============================================================
\section{Clase 1 -- Bandidos de $k$ brazos}
%% ============================================================

El problema de \textbf{bandidos de $k$ brazos} es la versi\'on m\'as simple de RL: hay un \'unico estado y $k$ acciones posibles. En cada paso el agente elige una acci\'on $a$ y recibe una recompensa. Cada elecci\'on es \textbf{independiente} de las anteriores (no hay estados que cambien).

\subsection{Valor de una acci\'on}
El \textbf{valor real} de una acci\'on es la recompensa esperada al tomarla:
\[
    q_*(a) \doteq \mathbb{E}[R_t \mid A_t = a].
\]
Como $q_*(a)$ es desconocido, el agente mantiene una \textbf{estimaci\'on} $Q_t(a)$. El m\'etodo m\'as simple (\emph{sample-average}) promedia las recompensas obtenidas al tomar $a$:
\[
    Q_t(a) = \frac{\text{suma de recompensas al tomar } a \text{ antes de } t}{\text{cantidad de veces que se tom\'o } a}.
\]

\subsection{Greedy, $\varepsilon$-greedy y exploraci\'on}
\begin{itemize}[nosep]
    \item \textbf{Greedy:} elegir siempre $A_t = \argmax_a Q_t(a)$. Explota, pero nunca explora, y puede quedar atrapado en una acci\'on sub\'optima.
    \item \textbf{$\varepsilon$-greedy:} con probabilidad $1-\varepsilon$ elige la acci\'on greedy, con probabilidad $\varepsilon$ elige una acci\'on al azar. La exploraci\'on queda \textbf{distribuida} a lo largo de todo el proceso.
\end{itemize}

\subsection{Actualizaci\'on incremental}
En vez de guardar todas las recompensas, se actualiza la estimaci\'on paso a paso:
\[
    Q_{n+1} = Q_n + \frac{1}{n}\,\big[R_n - Q_n\big].
\]
Es un caso de la regla general de aprendizaje:
\[
    \text{NuevaEst} \leftarrow \text{ViejaEst} + \text{TasaApr}\,\big[\text{Objetivo} - \text{ViejaEst}\big].
\]
Ventaja: memoria y c\'omputo $O(1)$ por paso, sin necesidad de almacenar toda la historia.

\subsection{Problemas no estacionarios}
Si $q_*(a)$ cambia con el tiempo, conviene usar un \textbf{paso constante} $\alpha \in (0,1]$ en lugar de $1/n$:
\[
    Q_{n+1} = Q_n + \alpha\,\big[R_n - Q_n\big].
\]
Esto da un \textbf{promedio con pesos exponencialmente decrecientes}: da m\'as peso a las recompensas recientes, permitiendo seguir cambios en el ambiente.

\subsection{Valores iniciales optimistas}
Se inicializa $Q_1(a)$ con un valor alto para toda $a$. Como toda recompensa real ``decepciona'' frente a la estimaci\'on inflada, el agente es empujado a probar las dem\'as acciones: exploraci\'on inicial sistem\'atica \emph{sin} aleatoriedad. Limitaci\'on: es un mecanismo transitorio, poco \'util en problemas no estacionarios.

%% ============================================================
\section{Clase 2 -- Procesos de Decisi\'on de Markov (MDP)}
%% ============================================================

Los \textbf{MDP} extienden los bandidos incorporando \textbf{estados} $s \in \mathcal{S}$: modelan la \textbf{toma de decisiones secuenciales}, donde cada decisi\'on afecta a las siguientes. Ahora se estima:
\begin{itemize}[nosep]
    \item $v_*(s)$: el \textbf{valor del estado} $s$.
    \item $q_*(s,a)$: el \textbf{valor de la acci\'on} $a$ en el estado $s$.
\end{itemize}

\subsection{Definici\'on formal}
Un MDP es una tupla $(\mathcal{S}, \mathcal{A}, \mathcal{R}, p)$, con estados, acciones y recompensas (conjuntos finitos), y $p$ la din\'amica. El tiempo es discreto: $t = 0,1,2,\dots$ Una \textbf{trayectoria} es una secuencia $S_0, A_0, R_1, S_1, A_1, R_2, \dots$

La din\'amica del ambiente queda caracterizada por:
\[
    p(s', r \mid s, a) \doteq \Pr\{S_t = s', R_t = r \mid S_{t-1} = s, A_{t-1} = a\},
\]
con $\sum_{s'}\sum_r p(s', r \mid s, a) = 1$. Si el ambiente es \textbf{determin\'istico}, esta probabilidad vale $1$ para un \'unico $(s', r)$; en general es \textbf{estoc\'astico}.

\subsection{Propiedad de Markov}
La probabilidad de $(S_t, R_t)$ depende \textbf{solo} de $(S_{t-1}, A_{t-1})$, no de la historia previa. Esto restringe c\'omo definimos los estados: deben incluir toda la informaci\'on del pasado relevante para el futuro.

\subsection{Retornos y episodios}
El objetivo del agente es maximizar el \textbf{retorno} (ganancia acumulada esperada):
\[
    G_t \doteq R_{t+1} + \gamma R_{t+2} + \gamma^2 R_{t+3} + \dots = \sum_{k=t+1}^{T} \gamma^{k-t-1} R_k,
\]
donde $\gamma \in [0,1]$ es el \textbf{factor de descuento}. Forma recursiva: $G_t = R_{t+1} + \gamma\, G_{t+1}$. En tareas \textbf{episódicas} hay un estado terminal ($T$ finito); en tareas \textbf{continuas} $T = \infty$.

\textbf{Factor de descuento $\gamma$:} es un par\'ametro del agente que fija cu\'anto valora recompensas cercanas vs.\ lejanas. Con $\gamma = 0$ el agente es \emph{miope} (solo cuenta la recompensa inmediata); con $\gamma \to 1$ valora m\'as las recompensas lejanas. \emph{Analog\'ia:} una tasa de inflaci\'on alta act\'ua como $\gamma$ bajo, licuando recompensas futuras.

\subsection{Pol\'iticas y funciones de valor}
Una \textbf{pol\'itica} $\pi(a \mid s)$ es la probabilidad de ejecutar $a$ en $s$ (puede ser determin\'istica o estoc\'astica).
\begin{align*}
    v_\pi(s) &\doteq \mathbb{E}_\pi[G_t \mid S_t = s] & \text{(valor del estado)} \\
    q_\pi(s,a) &\doteq \mathbb{E}_\pi[G_t \mid S_t = s, A_t = a] & \text{(valor de la acci\'on)}
\end{align*}
Propiedad: $q_\pi(s,a) = \mathbb{E}_\pi[R_{t+1} + \gamma\, v_\pi(S_{t+1}) \mid S_t = s, A_t = a]$.

\subsection{Pol\'iticas y valores \'optimos}
$\pi \geq \pi'$ sii $v_\pi(s) \geq v_{\pi'}(s)$ para todo $s$. Una pol\'itica $\pi_*$ es \textbf{\'optima} si domina a toda otra. Las \textbf{funciones de valor \'optimas} son:
\[
    v_*(s) \doteq \max_\pi v_\pi(s), \qquad q_*(s,a) \doteq \max_\pi q_\pi(s,a).
\]

\subsection{B\'usqueda de pol\'iticas \'optimas}
$v_*$ y $q_*$ son desconocidas; el agente las estima.
\begin{itemize}[nosep]
    \item \textbf{M\'etodos \emph{model-based}:} si conoci\'eramos $p(s', r \mid s, a)$ podr\'iamos usar programaci\'on din\'amica (\emph{Value Iteration}). Pero casi nunca conocemos $p$: interés te\'orico.
    \item \textbf{M\'etodos \emph{model-free}:} no necesitan $p$; aprenden interactuando. Ejemplos: Monte Carlo y Q-Learning.
\end{itemize}

%% ============================================================
\section{Clase 3 -- M\'etodos Monte Carlo}
%% ============================================================

Los m\'etodos \textbf{Monte Carlo} (MC) obtienen soluciones mediante \textbf{pruebas aleatorias repetidas}. En RL se basan en \textbf{generar episodios, observar y aprender de lo ocurrido}. Como no se conoce $p(s', r \mid s, a)$, el agente aprende con interacciones \textbf{reales} con el ambiente.

\subsection{Algoritmo MC para estimar $\pi_*$}
\begin{enumerate}[nosep]
    \item Inicializar $\pi$ al azar, $Q(s,a)$ arbitrario, y $\text{Retornos}(s,a)$ como listas vac\'ias.
    \item Generar un episodio completo siguiendo $\pi$: $S_0, A_0, R_1, \dots, S_{T-1}, A_{T-1}, R_T$.
    \item Para cada paso $t$: calcular $G_t = \sum_{k=t+1}^{T} \gamma^{k-t-1} R_k$, agregarlo a $\text{Retornos}(S_t, A_t)$, y actualizar $Q(S_t, A_t) \leftarrow \text{promedio}(\text{Retornos}(S_t, A_t))$.
    \item Mejorar la pol\'itica de forma greedy: $\pi(S_t) \leftarrow \argmax_a Q(S_t, a)$.
\end{enumerate}
$Q(s,a)$ converge a $q_*(s,a)$ a medida que la cantidad de visitas a $(s,a)$ tiende a infinito.

\subsection{Problema de exploraci\'on y versi\'on final}
Con una pol\'itica puramente greedy, muchos pares $(s,a)$ se visitar\'ian poqu\'isimo (o nunca), y sus estimaciones ser\'ian malas. Soluci\'on: usar una pol\'itica \textbf{$\varepsilon$-greedy} que garantice explorar todos los pares:
\[
    \pi(a \mid s) =
    \begin{cases}
        1 - \varepsilon + \dfrac{\varepsilon}{|\mathcal{A}(s)|} & \text{si } a = A^* \\[2mm]
        \dfrac{\varepsilon}{|\mathcal{A}(s)|} & \text{si } a \neq A^*
    \end{cases}
\]
donde $A^* = \argmax_a Q(s, a)$. La versi\'on final devuelve una pol\'itica estoc\'astica $\pi \approx \pi_*$.

%% ============================================================
\section{Clase 4 -- Diferencias Temporales: SARSA y Q-learning}
%% ============================================================

\subsection{MC vs.\ TD}
Ambos parten de la regla $\text{NuevaEst} \leftarrow \text{ViejaEst} + \text{TasaApr}\,[\text{Objetivo} - \text{ViejaEst}]$, pero difieren en \textbf{cu\'ando} y con qu\'e objetivo actualizan:
\begin{itemize}[nosep]
    \item \textbf{Monte Carlo:} espera al \textbf{final del episodio} y usa el retorno real $G_t$:
    \[
        Q(S_t, A_t) \leftarrow Q(S_t, A_t) + \alpha\,[G_t - Q(S_t, A_t)].
    \]
    \item \textbf{Diferencias Temporales (TD):} actualiza \textbf{tras cada paso}, sin esperar el final, usando $R_{t+1} + \gamma\, Q(S_{t+1}, A_{t+1})$ como objetivo:
    \[
        Q(S_t, A_t) \leftarrow Q(S_t, A_t) + \alpha\,[R_{t+1} + \gamma\, Q(S_{t+1}, A_{t+1}) - Q(S_t, A_t)].
    \]
\end{itemize}
\textbf{Bootstrapping:} estimar un valor a partir de \emph{otras estimaciones} propias ($Q(S_{t+1}, \cdot)$) en vez de esperar el resultado real. TD hace bootstrapping; MC no. Ventaja de TD: aprende ``sobre la marcha'', \'util en tareas continuas o de episodios muy largos.

\subsection{SARSA (on-policy)}
Estima $\pi_*$ usando \textbf{la misma pol\'itica} $\varepsilon$-greedy para elegir la acci\'on siguiente $A'$ y para estimar el retorno. La actualizaci\'on usa el par real $(S', A')$ que efectivamente se va a tomar:
\[
    Q(S, A) \leftarrow Q(S, A) + \alpha\,[R + \gamma\, Q(S', A') - Q(S, A)].
\]
Por usar la pol\'itica real en el objetivo, es un m\'etodo \textbf{on-policy}.

\subsection{Q-learning (off-policy)}
Cambia el objetivo por el m\'aximo sobre las acciones, desacoplando la estimaci\'on del retorno de la pol\'itica seguida:
\[
    Q(S, A) \leftarrow Q(S, A) + \alpha\,[R + \gamma \max_a Q(S', a) - Q(S, A)].
\]
Asume una pol\'itica greedy pura para estimar el retorno, sin importar qu\'e acci\'on se tome de verdad (que sigue siendo $\varepsilon$-greedy). Por eso es un m\'etodo \textbf{off-policy}. La tasa $\alpha \in (0,1]$ controla la relevancia de la informaci\'on reciente (\'util en ambientes no estacionarios).

%% ============================================================
\section{Clase 5 -- Deep Q-Learning (DQL)}
%% ============================================================

\subsection{L\'imite de los m\'etodos tabulares}
Los m\'etodos tabulares guardan un valor por cada par $(s,a)$ en una tabla $Q$. Esto es imposible cuando $\mathcal{S}$ es enorme (e.g.\ los p\'ixeles de una pantalla de Atari): la cantidad de estados es combinatoriamente gigantesca y, adem\'as, dos pantallas casi id\'enticas cuentan como estados distintos (no hay \textbf{generalizaci\'on}).

\subsection{Aproximaci\'on de la funci\'on de valor}
Se reemplaza la tabla por una funci\'on parametrizada $\hat q(s, a, \mathbf{w})$ (e.g.\ una red neuronal con pesos $\mathbf{w} \in \mathbb{R}^d$, con $d \ll |\mathcal{S}|$), que permite generalizar a estados nunca vistos. Se define el error cuadr\'atico:
\[
    \overline{\text{QE}}(\mathbf{w}) \doteq \sum_{s}\sum_{a} \big[q_*(s,a) - \hat q(s, a, \mathbf{w})\big]^2.
\]
Como $\hat q$ es diferenciable respecto de $\mathbf{w}$, se ajusta con \textbf{descenso por gradiente estoc\'astico} (SGD). Combinando SGD con bootstrapping ($q_*(S_t, A_t) \approx R_{t+1} + \gamma \max_a \hat q(S_{t+1}, a, \mathbf{w})$) se obtiene la regla \textbf{semi-gradiente}:
\[
    \mathbf{w} \leftarrow \mathbf{w} + \alpha\,[R + \gamma \max_a \hat q(S', a, \mathbf{w}) - \hat q(S, A, \mathbf{w})]\,\nabla \hat q(S, A, \mathbf{w}).
\]
No hay garant\'ia de convergencia (a diferencia del caso tabular), pero se usa mucho en la pr\'actica.

\subsection{Detalles de implementaci\'on (Mnih et al.\ 2013, 2015)}
El paper \textit{Playing Atari with Deep RL} / \textit{Human-level control} aprendi\'o a jugar Atari desde los p\'ixeles con una red convolucional. Dos t\'ecnicas clave para estabilizar el entrenamiento:
\begin{itemize}[nosep]
    \item \textbf{Replay de experiencias:} cada transici\'on $e_t = (S_t, A_t, R_{t+1}, S_{t+1})$ se guarda en una \emph{memoria} $D$; cada actualizaci\'on de $\mathbf{w}$ se hace sobre un \emph{minibatch} muestreado al azar. Esto rompe la correlaci\'on temporal entre muestras consecutivas.
    \item \textbf{Target network:} una copia de la red (\emph{Q-network}) usada para computar el objetivo, que se actualiza solo cada cierta cantidad de iteraciones. Mantener el objetivo ``fijo'' un rato evita perseguir un blanco que se mueve con cada paso.
\end{itemize}
El tama\~no de la memoria, el del minibatch y la frecuencia de actualizaci\'on de la target network son \textbf{hiperpar\'ametros} (adem\'as de $\alpha$, $\varepsilon$ y $\gamma$).

%% ============================================================
\section{Clase 6 -- Temas Avanzados}
%% ============================================================

\subsection{M\'etodos de gradiente de pol\'iticas}
En vez de aproximar $q_*(s,a)$ y derivar una pol\'itica, se aprende \textbf{directamente} una pol\'itica parametrizada $\pi(a \mid s, \boldsymbol{\theta})$. \'Util cuando hay demasiadas acciones o acciones \emph{continuas} (redes el\'ectricas, trading, rob\'otica, LLMs). Con una m\'etrica de desempe\~no $J(\boldsymbol{\theta})$ se hace ascenso por gradiente:
\[
    \boldsymbol{\theta}_{t+1} = \boldsymbol{\theta}_t + \alpha\, \nabla J(\boldsymbol{\theta}_t).
\]

\textbf{REINFORCE} (gradiente de pol\'iticas tipo MC): tras generar un episodio, para cada paso actualiza
\[
    \boldsymbol{\theta} \leftarrow \boldsymbol{\theta} + \alpha\, G\, \nabla \ln \pi(A_t \mid S_t, \boldsymbol{\theta}).
\]
Ajusta directamente la probabilidad de elegir $A_t$: sube o baja proporcionalmente al retorno observado $G$.

\textbf{Actor-Critic:} variante TD de REINFORCE. Reemplaza $G_t$ por el \textbf{error TD} usando bootstrapping:
\[
    \boldsymbol{\theta}_{t+1} \doteq \boldsymbol{\theta}_t + \alpha\,\big(R_{t+1} + \gamma\, \hat v(S_{t+1}, \mathbf{w}) - \hat v(S_t, \mathbf{w})\big)\, \nabla \ln \pi(A_t \mid S_t, \boldsymbol{\theta}).
\]
La red $\pi$ es el \textbf{actor} (elige acciones) y $\hat v$ el \textbf{cr\'itico} (eval\'ua estados). \textbf{PPO} (Proximal Policy Optimization) extiende Actor-Critic restringiendo cu\'anto puede cambiar $\pi$ por update: m\'as estable, estado del arte en muchos dominios.

\subsection{Monte Carlo Tree Search (MCTS)}
Para juegos como Go, la planificaci\'on cl\'asica (Minimax) no alcanza: $\sim$250 movidas por posici\'on y sin buenas heur\'isticas para podar. Idea: usar \textbf{trayectorias simuladas} desde el estado actual hasta el final, eligiendo acciones seg\'un una pol\'itica y aproximando el ambiente con un modelo (el agente \emph{imagina}, no interact\'ua con el ambiente real).

\textbf{M\'etodos Rollout:} simular trayectorias desde cada acci\'on de $S$, actualizar solo $(S,A)$, y luego elegir la acci\'on de mayor valor estimado. La \emph{pol\'itica rollout} es r\'apida (p.ej.\ aleatoria). No se almacenan las mejoras: no es aprendizaje real.

\textbf{MCTS} s\'i almacena las estimaciones MC en un \textbf{\'arbol} con ra\'iz en $S$, que se expande incrementalmente. Cuatro etapas por iteraci\'on:
\begin{enumerate}[nosep]
    \item \textbf{Selecci\'on:} bajar por el \'arbol siguiendo la \emph{pol\'itica del \'arbol} (balancea exploraci\'on/explotaci\'on, p.ej.\ $\varepsilon$-greedy) hasta una hoja.
    \item \textbf{Expansi\'on:} agregar uno o m\'as hijos a la hoja.
    \item \textbf{Simulaci\'on:} completar la trayectoria con la \emph{pol\'itica rollout}.
    \item \textbf{Backup:} propagar el retorno para actualizar los valores $(s,a)$ del \'arbol.
\end{enumerate}
Tras cierto tiempo, se elige en la ra\'iz la acci\'on de mayor valor. Parte del \'arbol se reusa; el resto se descarta.

\subsection{AlphaGo}
\begin{itemize}[nosep]
    \item \textbf{AlphaGo} (2015--2016): MCTS + redes neuronales profundas (\emph{rollout policy}, \emph{policy network}, \emph{value network}), entrenadas con 30M de partidas humanas + gradiente de pol\'iticas. Venci\'o a Fan Hui (5--0) y a Lee Sedol (4--1).
    \item \textbf{AlphaGo Zero} (2017): MCTS (sin rollouts) + una \'unica red (policy+value) entrenada por \textbf{self-play}, \emph{sin} dataset de partidas humanas. Aprendi\'o desde \textbf{tabula rasa}, venci\'o 100--0 a AlphaGo y descubri\'o estrategias nuevas.
\end{itemize}

\subsection{Preguntas para las Neurociencias}
El RL abre preguntas sobre el cerebro: \textquestiondown c\'omo eligen los animales entre opciones con recompensas desconocidas? \textquestiondown Hay un mecanismo parecido a $\varepsilon$-greedy? \textquestiondown C\'omo se relaciona la recompensa de RL con la \textbf{dopamina}? \textquestiondown Representa el cerebro algo como una tabla $Q$ o una funci\'on de valor? \textquestiondown Existe un an\'alogo al \emph{replay de experiencias} (consolidaci\'on de memoria durante el sue\~no, rol del hipocampo en la planificaci\'on con rollouts)? \textquestiondown Hay estructuras an\'alogas al actor y al cr\'itico?

\end{document}
